\begin{frame}{스무딩 효과 숫자로: `당첨'이 들어온 메일}
  앞의 학습 데이터(어휘 \{공짜, 회의\})에서 \textbf{한 번도 보지
  못한} `당첨'이 포함된 새 메일 $\{$공짜, \textbf{당첨}$\}$:
  \vskip0.3em
  \small
  \begin{tabular}{@{}lll@{}}
    \toprule
     & $P(\text{당첨 있음}|\text{spam})$ & $P(\text{당첨 있음}|\text{ham})$ \\
    \midrule
    \textbf{스무딩 없이} & $0/2 = \textbf{0}$ & $0/2 = \textbf{0}$ \; $\Rightarrow$ 0 vs 0 \textbf{비교 불가} \\
    \textbf{라플라스} & $1/4$ & $1/4$ \; $\Rightarrow$ \textbf{중립} 신호 \\
    \bottomrule
  \end{tabular}
  \vskip0.8em
  \begin{itemize}
    \item 스무딩을 쓰면 `당첨'은 \textbf{중립}이 되고, 남은 `공짜'
      (스팸 3/4 vs 정상 1/4)이 판단을 결정:
      $P(x|\text{spam}) = \tfrac34\times\tfrac14 = 0.1875$,
      $P(x|\text{ham}) = \tfrac14\times\tfrac14 = 0.0625$
      --- `공짜'가 본래 가질 \textbf{3배의 증거}(0.1875/0.0625)를
      그대로 살려 \textbf{스팸}으로 분류.
    \item ``0/0 비교 불가''에서 ``3:1로 스팸''으로 바뀌는 것 ---
      스무딩은 ``확률 0의 사건''을 ``아직 관찰되지 않았을 뿐, 매우
      희소할 것''으로 \textbf{재정의}.
  \end{itemize}
\end{frame}
